Zu jedem Punkt

\mavergleichskette
{\vergleichskette
{ P
}
{ \in }{ M
}
{  }{
}
{    }{
}
{   }{
}
}
{}{}{}
gibt es eine offene Kartenumgebung

\mavergleichskette
{\vergleichskette
{ P
}
{ \in }{ U
}
{  }{
}
{    }{
}
{   }{
}
}
{}{}{}
und eine Kartenabbildung
\maabbdisp {\alpha} { U } { V
} {}
mit

\mavergleichskette
{\vergleichskette
{ V
}
{ \subseteq }{ \R^n
}
{  }{
}
{    }{
}
{   }{
}
}
{}{}{}
offen und so, dass

\mavergleichskette
{\vergleichskette
{ \alpha^{-1 *} \omega
}
{ = }{ fdx_1 \wedge \ldots \wedge dx_n
}
{  }{
}
{    }{
}
{   }{
}
}
{}{}{}
ist mit
\mathl{f}{} stetig und positiv. Wir finden auch eine offene Umgebung

\mavergleichskette
{\vergleichskette
{ P
}
{ \in }{ U'
}
{ \subseteq }{ U
}
{    }{
}
{   }{
}
}
{}{}{,}
die homöomorph zu einem offenen Ball

\mavergleichskette
{\vergleichskette
{ U'
}
{ \cong }{ B'
}
{ \subseteq }{ V
}
{    }{
}
{   }{
}
}
{}{}{}
ist, wobei man auch annehmen kann, dass der Abschluss des Balles ganz in $V$ liegt. Der abgeschlossene Ball ist abgeschlossen und beschränkt, daher ist die stetige Funktion $f$ darauf und somit auch auf $U'$ beschränkt. Es folgt, dass
\mathl{\int_{U'} \omega}{} endlich ist, wobei $U'$ eine offene Umgebung von $P$ ist.

Diese offenen Mengen

\mavergleichskette
{\vergleichskette
{ U'
}
{ = }{ U'(P)
}
{  }{
}
{    }{
}
{   }{
}
}
{}{}{}
überdecken $M$. Wegen der Kompaktheit gibt es eine endliche Überdeckung

\mavergleichskettedisp
{\vergleichskette
{ M
}
{ =} { \bigcup_{i = 1} ^n U_i
}
{ } {
}
{ } {
}
{ } {
}
}
{}{}{}
mit

\mavergleichskettedisp
{\vergleichskette
{ \int_{U_i } \omega
}
{ <} { \infty
}
{ } {
}
{ } {
}
{ } {
}
}
{}{}{.}
Wegen der Positivität gilt somit

\mavergleichskettedisp
{\vergleichskette
{ \int_{M} \omega
}
{ \leq} { \sum_{i = 1}^n  \left(  \int_{U_i } \omega  \right)
}
{ <} { \infty
}
{ } {
}
{ } {
}
}
{}{}{.}

 [[Kategorie:Latexseite]]
